\documentclass[times,twocolumn,final]{elsarticle}

\usepackage{cag}
\usepackage{framed,multirow}
\usepackage{amssymb}
\usepackage{amsmath}
\usepackage{lineno}
\usepackage{graphicx}
\usepackage{booktabs}
\usepackage{placeins}

\newtheorem{proposition}{Proposition}
\newtheorem{remark}{Remark}

\journal{Computers \& Graphics}

\begin{document}

\begin{frontmatter}

\title{Adapter Scaling Trade-off and Timestep-Conditioned Scheduling in Multi-view Diffusion Texture Generation}

\author{}

\begin{abstract}
Multi-view diffusion models have become an important approach to 3D texture generation, where adapter modules are commonly used to inject geometric conditions such as normals and depth maps to improve cross-view consistency and geometric alignment. However, the adapter scaling strength during inference is often selected empirically, and its influence on shape fidelity and texture detail remains insufficiently analyzed. This paper systematically studies adapter scaling in geometry-controlled multi-view texture generation. We find that a larger adapter scale can improve structural metrics such as FG-SSIM, but it may also suppress the texture synthesis capability of the base diffusion model, leading to weakened color variation, flattened surfaces, and loss of high-frequency details. This observation indicates that adapter scaling is not merely a structure-enhancement parameter, but a key control variable that governs the balance between shape consistency and texture preservation. Rather than treating this as a generic guidance-scheduling problem, we identify an adapter-specific temporal conflict: geometric residuals injected into intermediate features are most useful for mid-stage structure refinement, whereas excessive residual injection in early or late stages can respectively disturb global formation or suppress texture synthesis. To mitigate this conflict, we propose Timestep-Conditioned Adapter Scaling (TCAS), a training-free inference strategy that concentrates strong geometric correction in the middle denoising stage and uses conservative scales in the early and late stages. The selected schedule C3 is chosen only on a 24-object probe set by analyzing 17 scaling variants and is then transferred unchanged to a 300-object validation set, where no parameter search is performed. Experiments show that TCAS maintains shape fidelity close to aggressive uniform scaling while achieving better signal fidelity and more stable texture preservation. These results suggest that, in adapter-augmented multi-view diffusion pipelines, adapter strength should be treated as a stage-dependent inference control variable rather than a fixed hyperparameter. TCAS requires no retraining, introduces negligible computational overhead, and can be used as a plug-and-play strategy for geometry-controlled multi-view texture generation tasks.
\end{abstract}

\begin{keyword}
\KWD
multi-view diffusion model \sep 3D texture generation \sep adapter scaling \sep timestep-conditioned scaling \sep shape-texture trade-off \sep geometric control \sep texture quality assessment
\end{keyword}

\end{frontmatter}

\linenumbers

\section{Introduction}
\label{sec:intro}

Recent advances in single-image 3D generation have substantially improved the recovery of geometry, multi-view consistency, and sparse-view reconstruction \cite{liu2023zero1to3,shi2023zero123pp,liu2023syncdreamer,long2023wonder3d,shi2024mvdream,wang2023imagedream,xu2024instantmesh,ke2024era3d}. As geometric generation models become increasingly mature, generating high-quality textures on an existing 3D shape has become a key factor affecting the realism and usability of 3D content. For an object whose geometry has already been determined, the final visual quality depends not only on whether the shape is complete, but also on whether the texture is faithful to the reference image, aligned with the 3D surface, and rich in local details.

Existing multi-view texture generation methods usually generate RGB images from multiple target views using 2D image diffusion models, and then map the generated images back to the 3D surface through projection or texture baking \cite{chen2023text2tex,richardson2023texture,poole2023dreamfusion,chen2023fantasia3d}. This route benefits from the generative capability of image diffusion models, but it still faces challenges in reference appearance alignment, geometry-texture consistency, and local texture detail preservation. Geometry-controlled multi-view diffusion methods represented by MVPainter introduce geometric conditions such as normals and depth maps to improve the alignment between generated textures and the underlying 3D surface, demonstrating the importance of explicit geometric control for 3D texture generation \cite{shao2025mvpainter}.

However, introducing geometric conditions is not only a question of whether a control signal should be used, but also a question of how strongly it should be applied. In adapter-augmented diffusion pipelines, geometric conditions are usually injected into the model through an adapter or control branch, and their influence is regulated by a scaling factor. In practice, this factor is often set empirically: a small scale may lead to insufficient geometric constraints, whereas a larger scale is usually expected to improve structural consistency.

In multi-view texture generation, however, stronger adapter scaling is not necessarily better. Texture generation requires not only contours and structures to be consistent with geometry, but also surface color variations, material details, and high-frequency patterns to be preserved. Overly strong geometric control may force the model to follow normal, depth, or edge conditions too aggressively, thereby suppressing the intrinsic texture synthesis capability of the base diffusion model. In this case, some structural metrics may still improve, while the visual result suffers from texture flattening, reduced color variation, and detail loss.

Therefore, adapter scaling in geometry-controlled multi-view diffusion deserves further analysis. Compared with prior work that mainly studies how to build texture generation pipelines, how to introduce geometric conditions, and how to improve overall generation quality, adapter strength is a fine-grained inference-time control problem. It directly affects the balance between shape consistency and texture detail, and is an easily overlooked factor in multi-view texture generation quality.

This problem is related to, but not solved by, existing guidance scheduling strategies. CFG scheduling changes the extrapolation between conditional and unconditional predictions, while a geometry adapter injects residuals into intermediate representations. Generic ControlNet or adapter scale scheduling similarly treats condition strength mainly as a condition-adherence knob. In texture generation, however, the key question is not only how strongly to follow the geometry, but also when geometric residuals should dominate and when they should recede so that the base model can synthesize color, material, and high-frequency texture. Simple monotonic schedules such as linear or cosine decay do not explicitly address this non-monotonic requirement.

The goal of this paper is not to propose a new full texture generation system, but to analyze and improve inference-time adapter scaling within a fixed geometry-controlled multi-view diffusion pipeline.

The main contributions of this paper are as follows:

(1) We systematically diagnose the shape-texture trade-off caused by adapter scaling in multi-view diffusion texture generation, showing that higher scales can improve structural metrics while simultaneously causing texture degradation whose specific form (flattening, artifact amplification, or no effect) depends on the adapter's residual magnitude and scale semantics.

(2) We identify an adapter-specific temporal conflict between geometric residual injection and texture formation, and propose Timestep-Conditioned Adapter Scaling (TCAS), which concentrates strong geometric correction in the middle denoising stage while using conservative scales in the early and late stages.

(3) Through uniform scale sweeps, texture audits, a 24-object probe analysis of 17 uniform, layer-wise, and temporal scaling variants, and a 300-object validation without re-searching the schedule, we show why fixed high scales and simple stage transfers of guidance scheduling are insufficient for the shape-texture trade-off.

(4) We further provide a large-scale texture preservation analysis on the 300-object validation set, showing that TCAS substantially reduces texture flattening compared with full high-scale scaling: C3 outperforms $s = 2.50$ on 87.7\%, 95.0\%, and 83.0\% of objects in Laplacian variance, RGB standard deviation, and gradient magnitude, respectively, while the 95\% confidence interval for the FG-SSIM difference crosses zero, indicating no significant difference in structural similarity.

\section{Related Work}
\label{sec:related}

\subsection{3D Texture Generation and Multi-view Diffusion}
\label{subsec:3dtex}

Single-image 3D generation is often decomposed into geometry generation and texture generation. In recent years, single-image-to-3D, multi-view diffusion, and sparse-view reconstruction methods have made notable progress in structural stability, shape completeness, and cross-category generalization \cite{liu2023zero1to3,shi2023zero123pp,liu2023syncdreamer,long2023wonder3d,shi2024mvdream,wang2023imagedream,xu2024instantmesh,ke2024era3d}. Nevertheless, high-quality texture generation remains a key bottleneck for the realism and usability of 3D content. Existing texture generation methods usually rely on 2D diffusion models to synthesize RGB images from multiple views, and then project or bake these images onto the 3D surface \cite{chen2023text2tex,richardson2023texture,poole2023dreamfusion,chen2023fantasia3d}. While this strategy inherits the strong generative capability of image diffusion models, it still suffers from reference appearance misalignment, geometry-texture inconsistency, and insufficient local texture detail.

To address these issues, methods such as MVPainter formulate 3D texture generation as a geometry-controlled multi-view diffusion task \cite{shao2025mvpainter}. MV-Adapter further shows that adapter-style multi-view control can enhance cross-view consistency without fully retraining the base model \cite{huang2024mvadapter}. By injecting geometric information such as normals and depth maps into the diffusion model as conditional signals, the generated results better follow the 3D surface structure and remain consistent across views. Unlike prior methods that focus mainly on multi-view image generation or general 3D reconstruction, these works emphasize precise alignment between texture and geometry as well as local detail recovery, significantly improving 3D texture generation quality. However, they mainly answer how to introduce geometric conditions and how to construct a texture generation pipeline, while the influence of geometric control strength at inference time on shape and texture quality remains underexplored.

\subsection{Adapter-augmented Diffusion Models and Geometric Control}
\label{subsec:adapter}

Adapter modules have become an important mechanism for introducing additional conditions into diffusion models. ControlNet injects spatial conditions such as edges, depth maps, and normals into a pretrained diffusion model through a parallel branch \cite{zhang2023controlnet}. IP-Adapter introduces image prompts through decoupled attention \cite{ye2023ipadapter}. LoRA and its variants implement lightweight model adaptation through low-rank residual updates \cite{hu2022lora,liu2024dora,zhang2023adalora}. Although these methods differ in architecture, they generally include some form of scaling factor that controls the influence of the external condition or adapter residual on the base model.

In practice, the adapter scale is often treated as an empirical hyperparameter. A larger scale usually improves condition adherence, for example by better following normals, depth maps, or reference images. In texture generation tasks, however, overly strong conditional constraints may weaken the base diffusion model's ability to synthesize color, material, and high-frequency details. Adapter scaling is therefore not a simple matter of increasing the strength as much as possible; it involves a balance between geometric constraints and texture fidelity. This paper focuses on exactly this inference-time control problem: given a fixed adapter structure and fixed geometric conditions, how should the adapter strength be set or scheduled to preserve both shape consistency and texture detail?

\subsection{Diffusion Guidance Scheduling and Timestep Control}
\label{subsec:guidance}

The denoising process of diffusion models is highly stage-dependent. Early steps mainly determine global layout and coarse structure, middle steps refine shape and semantic structure, and late steps are more responsible for color, texture, and high-frequency details \cite{ho2020ddpm,song2021ddim}. Based on this observation, many guidance scheduling strategies have been explored in both research and practice. For example, classifier-free guidance scales may be adjusted linearly, with cosine schedules, or dynamically to mitigate oversaturation, detail loss, or reduced diversity caused by overly strong guidance \cite{ho2022cfg,castillo2023adaptive}. In addition, the development of latent diffusion, high-resolution diffusion, and Transformer-based diffusion models further shows the close relationship among generation quality, noise schedules, conditional injection, and network architecture \cite{rombach2022ldm,esser2024scaling,blackforestlabs2024flux,loshchilov2019decoupled,cherti2023reproducible,nichol2022glide,saharia2022photorealistic,peebles2023scalable}.

These methods all exploit the staged nature of denoising to regulate conditional strength. However, different control signals act through different mechanisms. CFG scheduling adjusts the weight between conditional and unconditional predictions, mainly affecting the balance among semantic control, diversity, and visual fidelity. By contrast, the scale of a geometric adapter acts directly on intermediate features or residual branches, controlling how strongly geometric conditions intervene in the generation process. ControlNet-style scale control is closer in spirit, but it is still usually interpreted as a condition-adherence strength; directly applying a larger, smaller, or smoothly decayed control scale does not specify which denoising stage should receive strong geometric residuals for texture generation.

This distinction matters because the desired schedule for geometry-controlled texture generation is not necessarily monotonic. A monotonic decay or warm-up can reduce over-control in one part of the trajectory, but it cannot simultaneously avoid unstable early residuals, strengthen the middle stage where geometry refinement is most useful, and weaken late residuals that compete with color and high-frequency texture formation. Thus, timestep-conditioned scheduling of adapter strength cannot be treated as a direct transfer of CFG scheduling or generic ControlNet scale scheduling; it requires a dedicated analysis of how geometric residual injection interacts with texture synthesis.

\section{Method}
\label{sec:method}

\subsection{Task Definition and Base Pipeline}
\label{subsec:task}

This paper builds on the MVPainter-style geometry-controlled multi-view texture generation framework \cite{shao2025mvpainter} and studies how the inference-time strength of geometric condition injection affects the generated results. Given a single reference image and the corresponding 3D geometry, the base pipeline first generates RGB texture images from six target views and then obtains a textured mesh through projection or texture baking. We do not modify the base UNet, VAE, or pretrained adapter weights. Instead, the adapter scaling factor is treated as the key inference-time control variable.

\begin{figure*}[t]
\centering
\includegraphics[width=\textwidth]{figures/fig1.pdf}
\caption{Method overview. GeoTex-Adapter generates geometric residuals and injects them into UNet features. TCAS dynamically adjusts the residual strength according to denoising stages to balance geometric consistency and texture fidelity. The final multi-view outputs are projected and baked into a textured 3D mesh.}
\label{fig:overview}
\end{figure*}

During multi-view diffusion, the 3D geometry is rendered into normal maps, depth maps, and foreground masks, denoted as:
\begin{equation}
\mathbf{G} = \{G_\text{normal}, G_\text{depth}, G_\text{mask}\}.
\end{equation}

The geometric adapter maps these conditions to a residual correction term and injects it into the intermediate hidden states of the UNet. For the $t$-th denoising step, the adapter injection can be written as:
\begin{equation}
\mathbf{h}'_t = \mathbf{h}_t + s(p_t) \cdot A_\phi(\mathbf{h}_t, \mathbf{G}),
\label{eq:injection}
\end{equation}
where $\mathbf{h}_t$ and $\mathbf{h}'_t$ denote the hidden states before and after injection, $p_t$ denotes the sampling progress, and $s(p_t)$ is the adapter scaling strength at the current denoising stage. This formulation highlights a key difference from CFG-style guidance: changing $s(p_t)$ does not reweight two predicted scores, but changes the magnitude of a geometry-conditioned residual added to the feature stream. A larger $s$ usually enhances geometric condition adherence. However, if strong residuals are continuously applied in the late denoising stage, they may suppress the synthesis of color, material, and high-frequency details by the base diffusion model. This paper therefore studies how to schedule $s$ during denoising to balance shape consistency and texture detail preservation.

\subsection{Shape-Texture Diagnostic Metrics}
\label{subsec:metrics}

The influence of adapter scaling on generated results cannot be evaluated by a single structural metric. We divide evaluation metrics into three groups. The first group consists of shape and structure metrics, including FG-SSIM, Edge-SSIM, and Crop-SSIM, which measure consistency with ground-truth renderings in the foreground, edge regions, and local foreground crops. The second group contains signal fidelity and perceptual similarity metrics, including PSNR and LPIPS, which reflect pixel-level error and differences in deep feature space. The third group consists of texture diagnostic metrics, including Laplacian variance, RGB standard deviation, and gradient magnitude, which are used to detect weakened color variation, surface smoothing, and high-frequency detail loss.

Taking Laplacian variance as an example, the foreground texture response is defined as:
\begin{equation}
M_\text{tex} = \text{Var}(\text{Lap}(I_\text{out} \odot G_\text{mask})),
\end{equation}
where Lap denotes the Laplacian filter, $G_\text{mask}$ denotes the foreground mask, and Var denotes spatial variance. A larger $M_\text{tex}$ usually corresponds to richer local high-frequency variation. If structural metrics improve while Laplacian variance, RGB standard deviation, or gradient magnitude decreases, the improvement may mainly come from contour smoothing or boundary convergence rather than true texture quality improvement.

\subsection{Timestep-Conditioned Adapter Scaling}
\label{subsec:tcas}

The diffusion denoising process is stage-dependent: early steps mainly establish the global layout and coarse contours; middle steps refine geometric structure and cross-view alignment; late steps are more responsible for color, material, and local texture synthesis. Based on this observation, we propose Timestep-Conditioned Adapter Scaling (TCAS), which replaces a globally fixed adapter scale with a stage-dependent scale function. TCAS is not intended as a generic smooth guidance schedule. Its low-high-low form encodes the adapter-specific hypothesis that geometry residuals should be strongest only after coarse structure becomes stable and before late texture synthesis dominates.

To avoid ambiguity caused by different raw timestep directions in different schedulers, we use sampling progress $p$ to represent the denoising process, where $p = 0$ denotes the start from noise and $p = 1$ denotes the end of sampling close to the final image state. TCAS divides denoising into early, middle, and late stages:
\begin{equation}
s(p) = \begin{cases} s_e, & 0 \leq p < 1/3 \\ s_m, & 1/3 \leq p < 2/3 \\ s_l, & 2/3 \leq p \leq 1 \end{cases}.
\label{eq:tcas}
\end{equation}

The selected C3 configuration used in this paper is:
\begin{equation}
(s_e, s_m, s_l) = (1.25, 2.50, 1.25).
\end{equation}

This design follows three principles. First, the early stage uses a conservative scale to avoid disrupting global structure formation with strong geometric residuals. Second, the middle stage uses a higher scale to strengthen geometric constraints from normals, depth maps, and foreground contours. Third, the late stage reduces the scale again to prevent geometric residuals from suppressing color variation, material detail, and high-frequency texture synthesis. This non-monotonic schedule is therefore different from simply decaying or increasing the adapter strength across all timesteps. TCAS introduces no new parameters and does not modify the base model or adapter weights; it only replaces the scaling factor during inference, so its computational overhead is negligible.

\paragraph{Why mid-stage concentration is selected: the CAI model.}
The low--high--low form is not assumed but derived from measurable stage utilities; we state the argument formally. Let $Q(s_e,s_m,s_l)$ denote a fidelity objective (PSNR here), let $L$ denote the fixed-low schedule, and let $z_k\in\{0,1\}$ indicate whether the high scale is assigned to stage $k\in\{e,m,l\}$.

\begin{proposition}[Guardrailed bang--bang optimality]
\label{prop:cai}
Suppose (i) the stage-utility model $Q(s_e,s_m,s_l)=Q(L)+\sum_{k}U_k z_k$, where $U_k=Q(\text{high scale only in stage }k)-Q(L)$ is the measured finite-difference utility of stage $k$, and (ii) per-stage guardrails under which stage $k$ is admissible for the high scale only if its intervention does not reduce FG-SSIM by more than $\tau$ nor raise the artifact-sensitive Laplacian-variance diagnostic beyond a factor $\rho$ relative to $L$. Then the admissible schedule maximizing $Q$ is bang--bang: $z_k^{\star}=\mathbf{1}[\,U_k>0 \text{ and stage } k \text{ admissible}\,]$. Moreover, for a stage whose utility is positive but statistically indistinguishable from zero and whose structural effect is insignificant, assigning the conservative scale yields a statistically indistinguishable fidelity at a strictly lower texture-deviation risk, so the conservative assignment is the risk-dominant choice.
\end{proposition}

\begin{quote}\small
\textit{Proof.} Under (i), $Q$ decomposes into the constant $Q(L)$ plus one independent binary term $U_k z_k$ per stage. Condition (ii) restricts each feasible set to $\{0,1\}$ (removing at most the value $1$ for inadmissible stages) without coupling the stages; the joint maximizer over a product of independent binary feasible sets is therefore the product of per-stage maximizers, i.e., $z_k^{\star}=1$ exactly when $k$ is admissible and $U_k>0$. For the tie-break, if the CI of $U_k$ contains $0$ and the structural change is insignificant, both assignments give statistically indistinguishable $Q$, while $z_k=0$ strictly reduces the texture-deviation risk. \hfill$\square$
\end{quote}

\begin{remark}
The utilities are measured finite differences on a 24-object stage ablation of the strong-residual adapter instance of Section~\ref{subsec:adapterdep}: $U_e=-0.788$~dB (95\% CI $[-1.420,-0.156]$) with a guardrail violation (FG-SSIM $-0.105$, CI $[-0.141,-0.069]$; foreground Laplacian variance $3.85\times$ fixed-low), $U_m=+0.574$~dB $[0.391,0.756]$, and $U_l=+0.077$~dB $[0.006,0.149]$ with a structurally insignificant FG-SSIM change ($+0.001$, CI $[-0.002,0.003]$). Proposition~\ref{prop:cai} selects exactly the conservative--high--conservative assignment C3: the early stage is inadmissible, the middle stage is admissible with the dominant gain, and the late stage takes the conservative tie-break.
\end{remark}

\begin{remark}
Additivity is an approximation: stage effects interact, and the complete C3 schedule is therefore validated directly (Secs.~\ref{subsec:c3}--\ref{subsec:adapterdep}). Changing the adapter, its residual magnitude, or the scale semantics changes the measured $U_k$ and requires re-estimation. Proposition~\ref{prop:cai} thus provides a conditional optimality guarantee within the measured adapter regime---explaining why the mid-stage assignment follows from data rather than from a hand-tuned rule---and not a universal one.
\end{remark}

\subsection{GeoTex-Adapter and FAC Adaptive Extension}
\label{subsec:fac}

GeoTex-Adapter, TCAS, and FAC correspond to three complementary dimensions of adapter control. GeoTex-Adapter determines what geometric residual correction should be injected. TCAS determines when the residual should be strengthened or weakened along the denoising-time dimension. FAC further modulates the residual in a fine-grained manner along spatial positions and frequency components.

GeoTex-Adapter is used as a fixed geometric residual backbone. It receives a 5-channel geometric input consisting of a 3-channel normal map, a 1-channel depth map, and a 1-channel foreground mask. A multi-scale geometric encoder produces a feature pyramid, and residual corrections are injected into multiple resolution levels of the UNet. TCAS acts on the global temporal scale of these residuals to control the influence of geometric conditions at different denoising stages without changing the adapter architecture.

To further examine whether stage-wise adapter control can be extended to learned adaptive correction, we also conduct experiments with Full Adaptive Correction (FAC). FAC contains three lightweight modules: LTAG learns timestep-dependent gating, GSG generates spatial gates from geometric features, and FSC performs frequency-selective correction on adapter residuals. Unlike fixed TCAS, FAC finely modulates adapter residuals in the temporal, spatial, and frequency domains. In the extension experiments, FAC uses TCAS as the base temporal schedule and further adds GSG and FSC for spatial and frequency correction.

It should be emphasized that FAC requires learning a small number of additional parameters. Therefore, TCAS, which requires no training, remains the main method in this paper. FAC is used only as a learned adaptive extension; Section~\ref{subsec:facexp} reports a controlled re-examination that delimits when learned modulation can be expected to help.

\section{Experiments}
\label{sec:experiments}

\subsection{Experimental Setup}
\label{subsec:setup}

Experiments are conducted using an MVPainter-style geometry-controlled multi-view diffusion pipeline. Each object contains a single reference image and corresponding 3D geometry. During preprocessing, the 3D geometry is rendered into target-view normal maps, depth maps, foreground masks, and ground-truth RGB images. The model generates six target views at a time, and all comparisons with ground-truth renderings are performed under unified views, resolution, and foreground masks.

Except for the FAC extension experiment, all scaling experiments keep the object list, reference image, target views, random seed, base UNet, VAE, adapter weights, and number of denoising steps fixed; only the adapter scale or schedule is changed. The number of sampling steps is 50. C3 is selected using only a 24-object probe set; the same schedule is then frozen and applied to the 300-object validation set, where no parameter search is performed. This probe-to-validation protocol isolates the effect of adapter scaling from changes in model architecture, training data, sampling budget, or validation-set tuning.

Accordingly, the comparisons in this paper are comparisons among inference-time adapter scaling strategies under the same base pipeline, rather than a system-level benchmark among different texture generation architectures. Uniform adapter scales serve as the primary deployment baselines because they correspond to the most direct way of choosing a single fixed adapter strength at inference time. The probe variants further test whether the trade-off can be resolved by simple alternatives such as globally high scales, layer-wise redistribution, or moving high adapter strength to early, late, or multiple denoising stages. TCAS changes only how this strength is scheduled across denoising stages.

We use four evaluation sets: (1) a 50-object scale-sweep set for analyzing the overall trend of structural metrics under 13 uniform scales; (2) a 26-object texture audit set for comparing texture degradation under $s = 1.25$, $s = 2.25$, and $s = 2.50$; (3) a 24-object probe set for evaluating 17 uniform, layer-wise, and timestep-conditioned scheduling variants and selecting C3; and (4) a 300-object validation set for testing the selected C3 schedule, uniform baselines, and the FAC extension without re-searching TCAS.

Evaluation metrics are divided into three categories: shape and structure metrics (FG-SSIM, Edge-SSIM, and Crop-SSIM), signal and perceptual metrics (PSNR and LPIPS), and texture diagnostic metrics (Laplacian variance, RGB standard deviation, and gradient magnitude). All metrics are first computed over six target views and then averaged over views and objects. For texture-related conclusions, we report means, object-level win rates, and significance tests to avoid judging generation quality based on a single structural score.

The FAC extension uses the same object list and evaluation protocol as the 300-object validation, and its lightweight modules are trained only on the adapter's training pool, which is disjoint from all evaluation objects. TCAS is used as the main baseline, while LTAG, LTAG+GSG, and Full FAC are evaluated as ablation variants under the same paired protocol. The base model and original adapter remain frozen, and only lightweight adaptive correction parameters are learned.

\paragraph{Implementation and dataset details}
All objects are drawn from the Objaverse 3D asset collection: ground-truth multi-view images are rendered offline with Blender (Cycles engine, 17 views at $512\times512$, white background), and six target views (front, front-right, right, back, left, front-left) are used throughout the paper. The adapter's training pool contains 1{,}118 rendered objects and the evaluation pools described above are drawn independently; we verified that the training and evaluation identifier sets intersect in zero objects, so no evaluation object is seen during adapter training. Within the 300-object validation pool, the 24-object probe set corresponds to the first 24 identifiers ($\mathrm{obj}_{0000}$--$\mathrm{obj}_{0023}$) and the remaining 276 identifiers ($\mathrm{obj}_{0024}$--$\mathrm{obj}_{0299}$) form a holdout that is strictly disjoint from the probe; Section~\ref{subsec:adapterdep} reports transfer statistics on this disjoint holdout in addition to the full-pool tables. The 50-object scale-sweep set comprises the first 50 objects of the pool ($\mathrm{obj}_{0000}$--$\mathrm{obj}_{0049}$), and the 26-object audit set is a purposive subset of these 50 objects, selected to span the observed range of per-object $\Delta$FG-SSIM behaviour for manual inspection; all evaluation sets are subsets of the 300-object pool, and no training object appears in any of them. The base pipeline is an MVPainter-style geometry-controlled multi-view diffusion model built on a joint six-view latent UNet, used frozen together with its VAE. The GeoTex-Adapter takes a five-channel geometric input (3-channel normal map, 1-channel depth map, 1-channel foreground mask), produces a multi-scale residual pyramid, and injects corrections into multiple UNet resolution levels; for the adapter studied in the main tables, the requested scale is applied multiplicatively at every injected layer. Inference uses the Euler discrete scheduler with 50 denoising steps and a fixed random seed (42) shared by all schedules within an experiment, so all variants share the same initial latents; the model generates six views per object. FG-SSIM, Edge-SSIM (edge masks derived from depth/normal discontinuities), PSNR, and LPIPS (AlexNet backbone) are computed on foreground-masked regions after background normalization; the texture diagnostics (foreground Laplacian variance, RGB standard deviation, gradient magnitude) are computed on foreground pixels and compared against the same ground-truth renderings under identical views and masks.

\paragraph{Statistical reporting.}
Object-level paired comparisons report two-sided paired tests and 95\% confidence intervals constructed by percentile bootstrap over objects with 10{,}000 resamples. Where a mean is accompanied by $\pm$, the quantity is the standard deviation across objects. Preference-study intervals add a participant-level bootstrap level (Section~\ref{subsec:largescale}).

\subsection{Shape-Texture Trade-off of Adapter Scaling}
\label{subsec:tradeoff}

Table~\ref{tab:sweep} reports the results of a 13-scale uniform adapter sweep on 50 objects.

\begin{table}[t]
\centering
\caption{Uniform adapter scale sweep results.}
\label{tab:sweep}
\small
\resizebox{\columnwidth}{!}{%
\begin{tabular}{lcccc}
\toprule
Scale & $\Delta$FG-SSIM$\uparrow$ & POS & $\Delta$Crop-SSIM$\uparrow$ & $\Delta$PSNR$\uparrow$ \\
\midrule
0.75 & $-$0.013 & 20/50 & +0.113 & 7.53 \\
0.90 & +0.024 & 30/50 & +0.137 & 8.32 \\
1.00 & +0.047 & 34/50 & +0.147 & 8.79 \\
1.10 & +0.070 & 36/50 & +0.159 & 9.08 \\
1.15 & +0.071 & 36/50 & +0.157 & 9.25 \\
1.25 & +0.088 & 37/50 & +0.161 & 9.37 \\
1.35 & +0.100 & 39/50 & +0.163 & 9.63 \\
1.50 & +0.113 & 40/50 & +0.163 & 9.76 \\
1.60 & +0.120 & 42/50 & +0.166 & 9.83 \\
1.75 & +0.127 & 43/50 & +0.173 & 9.97 \\
2.00 & +0.134 & 44/50 & +0.169 & 9.96 \\
2.25 & +0.142 & 43/50 & +0.175 & 9.91 \\
2.50 & +0.145 & 44/50 & +0.165 & 9.91 \\
\bottomrule
\end{tabular}%
}
\end{table}

As the scale increases from $s = 0.75$ to $s = 2.50$, $\Delta$FG-SSIM rises from $-$0.013 to +0.145, and the number of objects with positive improvement increases from 20/50 to 44/50. This indicates that stronger adapter residuals indeed improve the model's adherence to geometric conditions such as normals, depth maps, and foreground contours. Figure~\ref{fig:curve} visualizes this trend: FG-SSIM keeps increasing with adapter scale, while Crop-SSIM saturates and slightly decreases at high scales.

\begin{figure}[t]
\centering
\includegraphics[width=\columnwidth]{figures/fig2.pdf}
\caption{Structural metric changes under uniform adapter scaling. The curves show that FG-SSIM keeps increasing with adapter scale, while Crop-SSIM saturates and slightly decreases at high scales. Stronger geometric control therefore mainly improves structural alignment, but does not necessarily keep improving local foreground quality.}
\label{fig:curve}
\end{figure}

Overall, the FG-SSIM improvement is stable, but the marginal benefit becomes smaller in the high-scale region. For example, increasing the scale from $s = 2.25$ to $s = 2.50$ yields only a limited FG-SSIM gain, while Crop-SSIM already decreases. This suggests that the nature of the high-scale benefit may change: the gain comes more from contour smoothing and foreground alignment than from preservation of local texture details. Figure~\ref{fig:heatmap} shows the per-object $\Delta$FG-SSIM heatmap, confirming that most objects gain higher structural scores as the scale increases, but with clear object-level differences.

\begin{figure}[t]
\centering
\includegraphics[width=\columnwidth]{figures/fig3.pdf}
\caption{$\Delta$FG-SSIM heatmap for 50 objects under 13 adapter scales. Most objects obtain higher structural gains as the scale increases, but there are clear object-level differences.}
\label{fig:heatmap}
\end{figure}

\subsection{Texture Audit: High Scale Causes Texture Flattening}
\label{subsec:audit}

To determine whether the FG-SSIM gain at high scales reflects better visual quality, we further conduct a texture audit on 26 objects, comparing texture metrics and foreground shape changes for $s = 2.50$ relative to $s = 1.25$. Figure~\ref{fig:qualitative} should be read as a local texture comparison rather than as a claim that any method fully matches the ground truth: the key evidence is the progressive loss of color variation, material marks, and high-frequency patterns inside the marked regions as the scale increases. The conservative scale $s = 1.25$ better preserves color variation and local texture, while $s = 2.50$ more easily produces uniform and smooth surfaces.

\begin{figure*}[t]
\centering
\includegraphics[width=\textwidth]{figures/fig4.pdf}
\caption{Texture degradation under different adapter scales. Each row shows one object; large crops highlight local texture regions (thumbnails below indicate crop locations). Green borders denote conservative scaling $s = 1.25$; red borders denote aggressive scaling $s = 2.50$. Within each row, increasing the adapter scale causes progressive loss of color variation and high-frequency detail: surfaces become more uniform and color-saturated, indicating that stronger geometric correction suppresses the base model's texture synthesis capability.}
\label{fig:qualitative}
\end{figure*}

The results show that 24 out of 26 objects (92\%) exhibit a pattern of ``texture loss without true shape gain.'' Only two objects show comparable texture quality, and no object shows a clear texture improvement from the high scale. Here, ``true shape gain'' is not equivalent to an increase in FG-SSIM. It refers to whether the high scale brings substantive improvement in effective foreground area, contour structure, or geometric detail compared with the conservative scale. If the metric improvement mainly comes from smoothed foreground boundaries or averaged contour errors without visible geometric improvement, it is categorized as having no true shape gain. Therefore, although a high adapter scale can improve structural metrics, such improvement does not necessarily indicate better visual shape or texture quality. Figure~\ref{fig:audit} further quantifies this texture loss.

\begin{figure}[t]
\centering
\includegraphics[width=\columnwidth]{figures/fig5.pdf}
\caption{Texture flattening at $s = 2.50$. Top: RGB standard deviation ratio on 26 objects ($s = 2.50$ / $s = 1.25$). Bottom: gradient magnitude ratio. Values below 1.0 indicate texture loss. Even the four best-improvement objects (green) show 24\% to 30\% texture loss despite positive FG-SSIM gains.}
\label{fig:audit}
\end{figure}

For multi-view texture generation, overly strong geometric control can suppress the color and detail synthesis capability of the base diffusion model during late denoising steps, making surfaces more uniform. Therefore, scale selection should not rely only on standard metrics such as FG-SSIM or PSNR, but should also incorporate texture-sensitive diagnostics.

\subsection{Timestep-Conditioned Adapter Scaling: Shape-Texture Balance of C3}
\label{subsec:c3}

Based on the texture audit, we do not pursue higher adapter scales throughout the entire denoising process. Instead, we further explore how different scaling strategies balance shape gain and texture loss. We construct 17 adapter scaling variants and evaluate them on the 24-object probe set. To avoid turning the 300-object validation into a hyperparameter search, the schedule selected in this probe stage is frozen before large-scale evaluation. These variants include three categories: uniform scaling baselines, layer-wise scaling strategies, and timestep-conditioned scaling strategies. Uniform scaling tests the effect of globally strengthening adapter influence and serves as the closest counterpart to using a single ControlNet-style condition scale. Layer-wise scaling examines whether redistributing the same geometric intervention across UNet levels is sufficient. Timestep scaling tests stage placement directly, including variants that put high adapter strength in early, middle, late, or broader temporal regions. This design is intended to rule out the simple explanation that any guidance-like schedule or any stronger control scale would solve the problem.

All variants are evaluated using $\Delta$FG-SSIM for shape gain and $\Delta$Lap Var for texture detail change. Delta texture metrics are computed relative to the corresponding ground-truth foreground renderings under the same views and masks. Importantly, the objective of this probe experiment is not to find the single highest structural score, but to identify a stable compromise strategy that avoids the trap of ``higher structural metrics but flattened texture.'' This criterion is crucial: if a schedule improves FG-SSIM by applying strong geometry residuals late but causes a clear drop in texture response, then it does not solve the adapter-specific conflict even though it appears favorable under a structure-only metric.

\begin{table}[t]
\centering
\caption{Shape-texture trade-off on the 24-object probe set. Delta texture metrics are computed relative to the corresponding ground-truth foreground renderings under the same views and masks; $\Delta$Lap Var closer to 0 indicates smaller texture loss.}
\label{tab:probe}
\resizebox{\columnwidth}{!}{%
\begin{tabular}{llccc}
\toprule
Variant & Category & $\Delta$FG-SSIM$\uparrow$ & $\Delta$Lap Var $\to$ 0 & Assessment \\
\midrule
C4 & Late-stage high & 0.105 & $-$0.0120 & Trap \\
C2 & Timestep sched. & 0.104 & $-$0.0110 & Trap \\
B2 & Layer-wise sched. & 0.090 & $-$0.0135 & Trap \\
A\_s2.00 & Uniform & 0.087 & $-$0.0062 & Trap \\
C3 (TCAS) & Mid-stage high & 0.084 & $-$0.0006 & Promising \\
A\_s2.25 & Uniform & 0.079 & $-$0.0096 & Trap \\
A\_s2.50 & Uniform & 0.078 & $-$0.0108 & Trap \\
B6 & Layer-wise sched. & 0.000 & $-$0.0027 & OK \\
\bottomrule
\end{tabular}%
}
\end{table}

Table~\ref{tab:probe} summarizes representative scheduling variants on the 24-object probe set. C4 and other late-stage high-scale strategies obtain higher $\Delta$FG-SSIM, but they are accompanied by obvious decreases in Laplacian variance, indicating that their structural gains come at the cost of texture loss. This is the main reason a direct transfer of guidance scheduling intuition is insufficient: a schedule can make the condition stronger at apparently useful timesteps and still damage the texture signal that matters for 3D appearance. Uniform high-scale strategies show a similar problem. For example, $s = 2.50$ achieves $\Delta$FG-SSIM of +0.078, but $\Delta$Lap Var decreases to $-$0.0108. In contrast, C3 achieves $\Delta$FG-SSIM of +0.084. Although it is slightly lower than some aggressive variants in structure gain, its Laplacian variance loss is only $-$0.0006, close to zero texture loss. Therefore, C3 is the only variant labeled Promising. We freeze this configuration before large-scale validation and do not tune it on the 300-object set.

The specific configuration of C3 uses the conservative scale $s = 1.25$ in the early and late stages, and the strong scale $s = 2.50$ in the middle geometry-refinement stage, i.e., C3 = (1.25, 2.50, 1.25). Its motivation follows the staged role of diffusion denoising. In the early stage, latent variables are still highly noisy, and strong geometric correction can interfere with global structure formation. In the middle stage, the signal becomes more stable, making it suitable to strengthen geometric constraints such as normals, depth maps, and foreground contours. In the late stage, generation mainly focuses on color, material, and fine texture synthesis, so adapter intervention should be reduced to preserve the texture generation space of the base diffusion model.

Figure~\ref{fig:c3compare} provides a qualitative comparison between C3 and uniform adapter scaling. The red-boxed zoom-in regions make the intended comparison visible: conservative scaling preserves more color variation but provides weaker structural correction, whereas aggressive scaling strengthens geometric constraints but flattens surface texture. C3 applies strong geometric correction in the middle stage and reduces adapter intervention in the late stage, thereby preserving more local texture than $s = 2.50$ while retaining stronger geometric correction than the conservative scale.

\begin{figure}[t]
\centering
\includegraphics[width=\columnwidth]{figures/fig6.pdf}
\caption{Qualitative comparison between C3 and uniform adapter scaling. Large crops show local texture regions (thumbnails below indicate crop locations). Columns: GT (gray), $s = 1.25$ (green), $s = 2.50$ (red), and C3 (orange). Compared with $s = 2.50$, C3 retains more local color variation and material texture while preserving stronger geometric correction than the conservative scale.}
\label{fig:c3compare}
\end{figure}

Figure~\ref{fig:pareto} further shows the distribution of the 17 variants in the shape-texture space. Most variants with high shape gain are accompanied by large texture loss and fall into the Trap region, whereas C3 maintains a relatively high shape gain while keeping texture loss near zero. Thus, the value of C3 is not that it maximizes every metric, but that it avoids the trap of ``higher structural scores but flattened texture,'' resulting in a more stable shape-texture compromise.

\begin{figure}[t]
\centering
\includegraphics[width=\columnwidth]{figures/fig7.pdf}
\caption{Shape-texture Pareto trade-off among 17 variants. The horizontal axis is shape gain, where farther right is better; therefore, the bottom-right region represents a more desirable compromise. C3 (star) lies in the Promising region, with high shape gain and nearly zero texture loss. Most high-shape-gain variants fall into the Trap region, where structural scores are obtained at the cost of large texture loss. Variants such as B6 have smaller texture degradation but almost no shape gain.}
\label{fig:pareto}
\end{figure}

Mechanistically, the advantage of C3 is not simply the result of hyperparameter tuning; it comes from the functional alignment between adapter correction and the diffusion process's stage-dependent objective. A natural hypothesis might be that correction should be strongest where the correction-to-noise ratio is highest. However, under the Euler sigma schedule, this ratio increases monotonically throughout denoising (from approximately 8 at full noise to over 400 near completion), which would predict that late-stage correction is most effective---a prediction inconsistent with the scheduling-variant results. The correct explanation involves \emph{functional alignment} rather than signal-to-noise ratio. The adapter correction is derived from geometric conditions (normals, depth, contours) and is therefore constructive only when the diffusion process is actively refining geometric structure. In the early stage ($\sigma > 5$), global structure has not yet emerged, so the correction has no stable target to reinforce and instead disrupts layout formation. In the middle stage ($\sigma \approx 5$--$1.2$), coarse structure is established and the denoiser is performing geometric refinement---the adapter correction direction aligns with this objective, making it maximally constructive. In the late stage ($\sigma < 1.2$), generation shifts to color, material, and fine texture synthesis; geometric correction is orthogonal to this objective and competes with the texture formation process. C3's low-high-low schedule therefore allocates strong correction proportional to functional alignment, not to signal clarity. The schedule is tied to the denoising process's functional role at each stage rather than to validation-set-specific tuning or a generic monotonic guidance rule.

This explanation is consistent with the scheduling-variant results. Early high-scale or early-and-late high-scale variants provide limited structural gain, indicating that early strong intervention is unstable. The late-stage high-scale configuration C4 obtains a high structural score but shows a significant decrease in Laplacian variance, making it a typical Trap variant where structural metrics improve while texture is flattened. Uniform high-scale variants similarly improve condition adherence but lose texture detail, showing that the issue is not solved by a stronger ControlNet-like scale. An aggressive C3 variant further increases the middle-stage strength but reduces signal fidelity, suggesting that even middle-stage correction has diminishing returns and can lead to over-constraint.

\FloatBarrier

\begin{table}[h!]
\centering
\caption{Comparison of representative variants under extended texture metrics. All delta texture metrics are computed relative to the corresponding ground-truth foreground renderings under the same views and masks. This table verifies that the advantage of C3 is not due to a single Laplacian variance metric.}
\label{tab:extended}
\resizebox{\columnwidth}{!}{%
\begin{tabular}{lcccc}
\toprule
Variant & $\Delta$Lap Var & $\Delta$RGB Std & $\Delta$Grad & $\Delta$HF Energy \\
\midrule
C3 & $-$0.001 & $-$0.023 & +0.142 & +0.058 \\
C4 & $-$0.012 & $-$0.032 & +0.086 & +0.062 \\
A\_s2.50 & $-$0.011 & $-$0.044 & +0.054 & +0.078 \\
A\_s2.00 & $-$0.006 & $-$0.040 & +0.100 & +0.070 \\
A\_s1.25 & $-$0.005 & $-$0.032 & +0.104 & +0.053 \\
\bottomrule
\end{tabular}%
}
\end{table}

C3 has nearly zero loss in Laplacian variance and remains relatively stable in RGB standard deviation, gradient magnitude, and high-frequency energy. Compared with C4 and uniform high-scale strategies, C3 shows weaker texture degradation, indicating that its shape-texture balance is not an artifact of a single metric.

\FloatBarrier

\subsection{Large-scale Validation}
\label{subsec:largescale}

Finally, as a transfer test rather than another search, we compare conservative uniform scaling ($s = 1.25$), aggressive uniform scaling ($s = 2.50$), and selected C3 (TCAS) on 300 objects. C3 is not re-searched on the 300-object set; it directly uses the timestep schedule selected on the 24-object probe set to evaluate the transferability of the strategy to a larger object set.

\begin{table}[t]
\centering
\caption{Large-scale validation on 300 objects. Values are absolute metrics of generated outputs with respect to ground truth.}
\label{tab:300obj}
\resizebox{\columnwidth}{!}{%
\begin{tabular}{lcccc}
\toprule
Method & FG-SSIM$\uparrow$ & Edge-SSIM$\uparrow$ & PSNR$\uparrow$ & FG-LPIPS$\downarrow$ \\
\midrule
$s = 1.25$ & 0.430 & 0.513 & 17.95 & 0.174 \\
$s = 2.50$ & 0.473 & 0.559 & 17.90 & 0.172 \\
C3 (TCAS) & 0.476 & 0.529 & 18.86 & 0.171 \\
\bottomrule
\end{tabular}%
}
\end{table}


As shown in Table~\ref{tab:300obj}, conservative uniform scaling $s = 1.25$ has lower FG-SSIM and Edge-SSIM, indicating insufficient geometric constraints. Aggressive uniform scaling $s = 2.50$ obtains the highest Edge-SSIM, suggesting that continuously strengthening adapter residuals helps edge and contour alignment. However, its PSNR is lower than that of C3, indicating that overly strong geometric intervention reduces overall signal fidelity. In contrast, C3 achieves an FG-SSIM of 0.476, which is comparable to $s = 2.50$, while increasing PSNR to 18.86 and slightly reducing FG-LPIPS. These results show that timestep-conditioned scaling preserves structural alignment while improving signal quality.

PSNR mainly reflects global pixel error and should not be directly equated with texture detail quality. To further verify texture preservation, we compute Laplacian variance, RGB standard deviation, and gradient magnitude on the same 300 objects to measure foreground high-frequency response, color variation amplitude, and local gradient strength, respectively. All results use the same object list, random seed, 50 denoising steps, and adapter weights to ensure comparability.


\begin{table}[t]
\centering
\caption{Large-scale texture preservation validation on 300 objects. Higher texture metrics indicate better preservation of foreground color variation and high-frequency details. Absolute values are means over objects and views; the ratios in Table~\ref{tab:retention} are the mean over objects of each per-object ratio relative to the corresponding $s = 1.25$ statistic (equal object weighting).}
\label{tab:texture300}
\small
\begin{tabular}{lccc}
\toprule
Method & Lap Var$\uparrow$ & RGB Std$\uparrow$ & Grad Mag$\uparrow$ \\
\midrule
$s = 1.25$ & 0.0151 & 0.0408 & 0.3258 \\
$s = 2.50$ & 0.0095 & 0.0251 & 0.2616 \\
C3 (TCAS) & 0.0129 & 0.0377 & 0.2919 \\
\bottomrule
\end{tabular}
\end{table}

Table~\ref{tab:texture300} shows that although $s = 2.50$ achieves the highest Edge-SSIM, it has the lowest values across all three texture metrics, indicating that full high-scale scaling weakens color variation and high-frequency details. C3 has FG-SSIM comparable to $s = 2.50$ while substantially improving Laplacian variance, RGB standard deviation, and gradient magnitude, demonstrating a better balance between structural alignment and texture preservation.


\begin{table}[t]
\centering
\caption{Texture retention ratios and object-level stability of C3. Each ratio is the mean over objects of the per-object ratio of the corresponding foreground statistic to that of $s = 1.25$ (equal object weighting); object-level win rates and confidence intervals for C3 over $s = 2.50$ are reported in the text.}
\label{tab:retention}
\resizebox{\columnwidth}{!}{%
\begin{tabular}{lccc}
\toprule
Method & Lap Var Ratio$\uparrow$ & RGB Std Ratio$\uparrow$ & Grad Mag Ratio$\uparrow$ \\
\midrule
$s = 1.25$ & 1.000 & 1.000 & 1.000 \\
$s = 2.50$ & 0.678 & 0.643 & 0.807 \\
C3 (TCAS) & 0.835 & 0.970 & 0.883 \\
\bottomrule
\end{tabular}%
}
\end{table}


As shown in Table~\ref{tab:retention}, compared with conservative scaling $s = 1.25$, aggressive scaling $s = 2.50$ retains only 67.8\% of Laplacian variance, 64.3\% of RGB standard deviation, and 80.7\% of gradient magnitude, indicating that full high-scale scaling substantially weakens foreground color variation and high-frequency details. In contrast, C3 retains 83.5\%, 97.0\%, and 88.3\% of these metrics, respectively, and clearly outperforms $s = 2.50$ across all three texture measures. This shows that TCAS still incurs some texture reduction, but it substantially alleviates texture flattening compared with full high-scale scaling.

Object-level statistics further show that C3 outperforms $s = 2.50$ on 87.7\%, 95.0\%, and 83.0\% of objects in Laplacian variance, RGB standard deviation, and gradient magnitude, respectively. The corresponding bootstrap 95\% confidence intervals do not cross zero, indicating statistically stable texture-preservation advantages. C3 also outperforms $s = 2.50$ on 74.3\% of objects in PSNR. By contrast, the confidence interval for FG-SSIM crosses zero, indicating no significant difference in foreground structural similarity between C3 and $s = 2.50$. Therefore, TCAS does not noticeably sacrifice structural alignment, while significantly reducing the texture degradation caused by full high-scale scaling.

In summary, the 300-object experiments demonstrate that adapter scaling affects both structural alignment and texture detail preservation. Full high-scale scaling strengthens edge constraints but is more likely to cause texture flattening. TCAS strengthens geometric correction in the middle stage and reduces adapter intervention in the late stage, thereby maintaining structural alignment while substantially reducing texture loss. It should be noted that C3 still has some texture decrease relative to $s = 1.25$; therefore, we do not claim that TCAS completely eliminates texture loss, but rather that it significantly reduces texture degradation compared with full high-scale scaling.

\paragraph{Perceptual validation with CLIP-IQA}
To address the concern that proxy texture statistics do not directly measure perceptual quality, we additionally evaluate CLIP-IQA, which builds on CLIP representations, on a 50-object validation subset as a complementary no-reference perceptual signal. CLIP-IQA is used as complementary perceptual evidence rather than a complete replacement for human preference evaluation.

\begin{table}[t]
\centering
\caption{CLIP-IQA perceptual validation on 50 validation objects. Scores are reported for foreground regions and full rendered views; $\pm$ denotes the standard deviation across the 50 objects. Higher is better.}
\label{tab:clipiqa}
\resizebox{\columnwidth}{!}{%
\begin{tabular}{lccc}
\toprule
Method & CLIP-IQA (FG)$\uparrow$ & CLIP-IQA (Full)$\uparrow$ & Win vs. $s = 2.50$ \\
\midrule
$s = 1.25$ & $\mathbf{0.4920 \pm 0.0032}$ & $\mathbf{0.4918 \pm 0.0030}$ & $\mathbf{98.0\%}$ \\
$s = 2.50$ & $0.4860 \pm 0.0018$ & $0.4857 \pm 0.0017$ & -- \\
C3 (TCAS) & $0.4906 \pm 0.0028$ & $0.4903 \pm 0.0025$ & 94.0\% \\
\bottomrule
\end{tabular}%
}
\end{table}

Table~\ref{tab:clipiqa} shows that C3 achieves a foreground CLIP-IQA score of 0.4906 versus 0.4860 for $s = 2.50$, scoring higher on 47 of the 50 objects (94.0\%). The paired comparison uses a two-sided paired $t$-test over objects ($t(49)=10.52$, $p\approx3.6\times10^{-14}$), and the 95\% confidence interval of the paired mean difference is constructed by object-level percentile bootstrap with 10{,}000 resamples, giving $[0.0037,\,0.0054]$, which excludes zero. The ordering $s = 1.25$ $>$ C3 $>$ $s = 2.50$ supports the diagnosis that aggressive scaling reduces perceptual quality, while TCAS recovers most of this loss with stronger geometry than conservative scaling.

\paragraph{Blinded preference study}
To further examine whether the diagnostic texture differences are perceptually meaningful, we conducted a blinded three-alternative forced-choice (3AFC) preference study on 30 validation objects. Twenty-four participants viewed, for each object and each of the three criteria (texture naturalness, shape consistency, overall quality), the anonymized outputs of conservative uniform scaling, aggressive uniform scaling, and C3 side by side; the left-to-right order of the three conditions was randomized independently per trial, and participants were not informed which method produced which image or how many conditions existed. Instructions asked participants to select the best image for the stated criterion only. Each participant judged all 30 objects for all three criteria, yielding 720 valid first-choice decisions for each criterion. The study complements the proxy diagnostics by testing whether the observed texture and structure differences are reflected in human preference.

\begin{table}[!htbp]
\centering
\caption{Blinded three-alternative first-choice preference study. Percentages are computed over 720 valid decisions per criterion.}
\label{tab:preference}
\resizebox{\columnwidth}{!}{%
\begin{tabular}{lccc}
\toprule
Method & Texture naturalness$\uparrow$ & Shape consistency$\uparrow$ & Overall quality$\uparrow$ \\
\midrule
$s = 1.25$ & $\mathbf{332/720\ (46.1\%)}$ & 96/720 (13.3\%) & 176/720 (24.4\%) \\
$s = 2.50$ & 104/720 (14.4\%) & $\mathbf{331/720\ (46.0\%)}$ & 126/720 (17.5\%) \\
C3 (TCAS) & 284/720 (39.4\%) & 293/720 (40.7\%) & $\mathbf{418/720\ (58.1\%)}$ \\
\bottomrule
\end{tabular}%
}
\end{table}
\FloatBarrier

The preference pattern in Table~\ref{tab:preference} follows the intended trade-off. Conservative scaling is most often selected for texture naturalness, while aggressive scaling is most often selected for shape consistency. C3 stays close to both single-scale baselines on their respective strengths and receives the highest overall preference, with 58.1\% of all overall-quality votes. Thus, TCAS is not claimed to maximize every individual criterion; it provides a perceptually preferred balance between geometric adherence and texture preservation.

Because first-choice decisions are clustered within both participants and objects, we additionally report 95\% confidence intervals obtained by a two-level cluster bootstrap (resampling participants and objects with replacement, 10{,}000 resamples). C3's overall-quality share is 58.0\% (95\% CI $[50.1\%, 65.8\%]$), exceeding the conservative and aggressive baselines by $+33.6$ (CI $[19.3, 47.2]$) and $+40.5$ percentage points (CI $[29.4, 51.8]$); both differences exclude zero, so the overall-preference advantage of C3 is significant under cluster-aware inference. By contrast, the baselines' single-criterion advantages are not significant (texture naturalness: $s = 1.25$ over C3 $+6.7$ points, CI $[-7.9, 21.5]$; shape consistency: $s = 2.50$ over C3 $+5.3$ points, CI $[-6.7, 17.9]$). Overall quality is therefore the only criterion with a significant preference, supporting the interpretation of TCAS as the perceptually preferred balance rather than a per-criterion maximum.

\subsection{FAC Adaptive Correction Extension: A Controlled Re-examination}
\label{subsec:facexp}

TCAS uses a fixed three-stage scale, which has the advantages of requiring no training, being simple to implement, and introducing low computational overhead. However, its residual scaling is still uniform across spatial positions and frequency components, lacking fine-grained adaptivity. To examine whether stage-wise adapter control can be extended to learned adaptive correction, we trained the FAC extension---LTAG for timestep gating, GSG for spatial gating, and FSC for frequency-selective correction on top of the TCAS temporal schedule. The base model, VAE, and GeoTex-Adapter remain frozen; only the three lightweight controllers are trainable, adding roughly $6\times10^{3}$ parameters. The controllers are trained with the same enhanced denoising objective as the base adapter (Adam, 2{,}000 steps) on the adapter's training pool, which is strictly disjoint from all evaluation objects, and all variants are evaluated on the strong-residual adapter instance of Section~\ref{subsec:adapterdep} under a strictly paired protocol (identical objects, seeds, and evaluation code).

In this controlled setting, no learned variant reached the training-free schedule. We tested seven configurations spanning the most favourable training conditions we could construct: joint training of all controllers; a two-stage scheme with the temporal gate warm-started at the exact C3 envelope and frozen while the spatial and frequency controllers train; an envelope penalty that constrains the gate to remain near C3 during joint training; a fidelity-weighted objective; and a gently scheduled variant with fourfold-reduced training. Their paired deficits relative to the TCAS baseline range from $-1.1$ to $-2.9$~dB PSNR (win rates at most 16\%), increasing monotonically with the amount of controller training. A warm-start control in which the gate is left untrained at the exact C3 envelope is statistically indistinguishable from TCAS (paired $\Delta=-0.18$~dB on the 12-object control, 95\% CI $[-0.42,+0.05]$, crossing zero), validating the evaluation path. The complete dose--response study is provided in the supplementary material.

The outcome delimits the scenario in which learned modulation affects the paper's conclusions. Preserving the temporal envelope is necessary but not sufficient: with the envelope preserved and the controllers untrained, the extension neither helps nor harms, while every trained configuration degraded fidelity, indicating that the remaining damage is driven by the learned-modulation training signal itself, which trades foreground fidelity for denoising-objective loss. Learning a modulation that improves upon the C3 schedule therefore requires a training signal that does not trade away foreground fidelity (e.g., fidelity-weighted objectives, envelope-preserving constraints, or post-hoc rather than gradient-learned gates), which we identify as future work. With the tested lightweight controllers, the fixed schedule is the stronger default, and TCAS, which is training-free and plug-and-play, remains the main and sole method of this paper.

\FloatBarrier

\subsection{Adapter-Dependence of High-Scale Failure Modes}
\label{subsec:adapterdep}

The preceding experiments demonstrate the texture-flattening failure mode under high adapter scaling on the studied adapter checkpoint. A natural question is whether this failure mode is universal or specific to the adapter's trained residual properties. We conduct additional analysis across adapter checkpoints with different residual magnitudes to characterize how the failure mode depends on adapter properties.

We find that the specific failure mode of aggressive adapter scaling is not universal. Across three adapter checkpoints spanning a spectrum of trained residual magnitudes, uniform high scale ($s = 2.50$) produced: (a) texture flattening when all layers are scaled without per-layer caps; (b) high-frequency artifact amplification when per-layer caps restrict the fine-texture layer to a low scale while coarse layers are boosted; or (c) no appreciable effect when the adapter's residual magnitude is very small. A per-layer scale-cap ablation on the same checkpoint reproduced both (a) and (b), isolating the inference-time scale semantics as the determining factor for the failure \emph{direction}.

A residual-magnitude sweep further shows that the intervention \emph{strength} is monotonic in the trained residual magnitude: scaling the adapter's output projection from 0.1$\times$ to 3$\times$ transitions the behavior from no-effect, to artifact amplification under high scale, to over-constraint even under conservative scale.

Moreover, a stage ablation shows that even \emph{which denoising stage is most sensitive to high scale} is adapter-dependent: on the artifact-prone adapter, early-stage high scale reproduced nearly all of the damage (FG-SSIM $\downarrow$ from 0.266 to 0.187) while late-stage high scale was nearly harmless (FG-SSIM 0.261), whereas on the paper's main adapter the late stage was the primary source of texture loss.

Despite this variability, in every regime where scaling had a non-negligible effect, the temporal low-high-low schedule C3 yielded the best shape-texture balance. This indicates that TCAS's value is structural rather than incidental: by relying on the diffusion process's stage-dependent roles (global layout $\to$ geometry refinement $\to$ texture synthesis), it mitigates the harm of high-scale intervention regardless of the adapter's specific failure mode. The failure mode (flatten vs.\ artifact vs.\ no-op), its strength, and its temporal location are all adapter-dependent; but the middle denoising stage being the optimal window for geometric correction appears to be a property of the diffusion process itself.

Two further robustness checks were performed on the strong-residual, per-layer-capped adapter instance described above. First, a boundary/peak sensitivity sweep (three high-scale windows $[0.25,0.75)$, $[1/3,2/3)$, $[0.4,0.6)$ $\times$ five peak scales from 2.00 to 3.00, fifteen candidates on the 24-object probe protocol) shows that the canonical middle-third window lies in the stable high-performance region: its 2.50 setting achieves 16.71~dB PSNR, the best setting in the same window (2.75) achieves 16.79~dB with worse texture diagnostics and lower FG-SSIM, and the broader and narrower windows reach at most 16.73 and 16.62~dB. C3 is therefore a robust middle-stage choice rather than a boundary-sensitive point estimate. Second, because the 24-object probe is part of the 300-object pool, transfer claims were re-verified on the disjoint 276-object holdout ($\mathrm{obj}_{0024}$--$\mathrm{obj}_{0299}$): C3 improves PSNR over the strongest non-C3 probe candidates by $+0.258$~dB (95\% CI $[0.207,0.308]$, 206/276 wins) over a trapezoid schedule and $+0.324$~dB ($[0.281,0.367]$, 232/276 wins) over a Gaussian-peak schedule, and over the reviewer-relevant generic schedules by $+0.323$~dB ($[0.263,0.386]$, 214/276 wins, FG-SSIM statistically tied) over linear warm-up and $+0.902$~dB ($[0.824,0.981]$, 253/276 wins, FG-SSIM $+0.046$) over the cosine bump. The holdout ordering matches the probe: warm-up approaches C3 in structure but loses in signal fidelity, and the smooth bump loses in both because its effective high-scale duration is shorter than the full middle third.


\section{Limitations}
\label{sec:limitations}

This paper mainly focuses on inference-time adapter scaling and does not retrain the base multi-view diffusion model or the adapter. Therefore, TCAS can alleviate texture degradation under high scales---which may manifest as flattening or artifact amplification depending on the adapter's residual properties---but it cannot fundamentally solve the limitations of the base model in invisible-region completion, complex material synthesis, and cross-view detail consistency.

Second, Laplacian variance, RGB standard deviation, and gradient magnitude are proxy metrics for texture degradation. They are effective for detecting high-frequency loss and surface smoothing, but they cannot fully replace human perceptual evaluation. For objects with strong reflection, transparency, thin structures, or highly stylized textures, discrepancies may exist between these metrics and subjective quality.

In addition, the current TCAS uses fixed stage boundaries and a fixed scale combination, C3 = (1.25, 2.50, 1.25). This setting is simple and stable, but it should be understood as a validated inference strategy for the studied geometry-adapter pipeline rather than a universally optimal schedule. As shown in Section~\ref{subsec:adapterdep}, the failure mode of high-scale intervention is adapter-dependent, and even the most sensitive denoising stage differs across adapters; yet C3's middle-stage concentration remains optimal in all tested regimes, suggesting that its temporal structure is robust even when specific failure mechanisms vary. It has not yet been adapted to object complexity, material type, input reference image quality, or different sampler families. Future work could study content-adaptive stage boundaries, SNR-aware scaling, and learned scale selection strategies.

Finally, our controlled re-examination of the learned FAC extension indicates that unconstrained lightweight gating can flatten the temporal envelope that makes TCAS effective; learned modulation is therefore promising only under envelope-preserving constraints, which we identify as a concrete direction for future work, while TCAS remains the main training-free method.


\section{Conclusion}
\label{sec:conclusion}

This paper presents a systematic experimental analysis of adapter scaling in multi-view diffusion texture generation. The experiments show that higher adapter scales continuously improve structural metrics such as FG-SSIM, but such improvement can mask texture quality degradation. In the 26-object texture audit, 92\% of objects show texture flattening and high-frequency detail loss under $s = 2.50$. This reveals an adapter-specific temporal conflict that is not captured by simply treating the scale as a generic guidance strength.

To address this shape-texture trade-off, we propose Timestep-Conditioned Adapter Scaling (TCAS). The selected configuration, C3, applies strong adapter correction in the middle denoising stage and uses conservative scales in the early and late stages. The 24-object probe set is used only to select C3 from 17 scaling variants that include uniform, layer-wise, and temporal alternatives. With this schedule fixed, the 300-object validation further shows that C3 achieves FG-SSIM comparable to aggressive uniform scaling $s = 2.50$ while improving PSNR by 0.96 dB.

Therefore, the core contribution of this paper is not to claim that the generated results are visually close to the ground truth in all cases, nor to present TCAS as a universal replacement for all guidance schedules. Instead, the contribution is to reveal and validate a hidden shape-texture trade-off caused by geometry-adapter residual injection, and to show that a simple non-monotonic inference schedule can mitigate this trade-off better than fixed high-scale control and other tested scheduling variants. This conclusion suggests that, in geometry-controlled multi-view diffusion models, adapter strength should be treated as a stage-dependent inference control variable rather than a single fixed hyperparameter.

As an additional check, a controlled re-examination of the learned-correction extension (FAC) under a strictly paired protocol showed that unconstrained lightweight gating does not preserve the schedule's shape; the analysis delimits the scenario in which learned modulation helps---only under envelope-preserving constraints---and the training-free TCAS remains the main and sole method.


\begin{thebibliography}{35}

\bibitem{liu2023zero1to3}
Liu R, Wu R, Van Hoorick B, et al. Zero-1-to-3: Zero-shot one image to 3D object. \textit{Proceedings of the IEEE/CVF International Conference on Computer Vision}, 2023.

\bibitem{shi2023zero123pp}
Shi R, Chen H, Zhang Z, et al. Zero123++: A single image to consistent multi-view diffusion base model. arXiv:2310.15110, 2023.

\bibitem{liu2023syncdreamer}
Liu Y, Lin C, Zeng Z, et al. SyncDreamer: Generating multiview-consistent images from a single-view image. arXiv:2309.03453, 2023.

\bibitem{long2023wonder3d}
Long X, Guo Y C, Lin C, et al. Wonder3D: Single image to 3D using cross-domain diffusion. arXiv:2310.15008, 2023.

\bibitem{shi2024mvdream}
Shi Y, Wang P, Ye J, et al. MVDream: Multi-view diffusion for 3D generation. \textit{International Conference on Learning Representations}, 2024.

\bibitem{wang2023imagedream}
Wang P, et al. ImageDream: Image-prompt multi-view diffusion for 3D generation. arXiv:2312.02201, 2023.

\bibitem{xu2024instantmesh}
Xu J, et al. InstantMesh: Efficient 3D mesh generation from a single image with sparse-view large reconstruction models. \textit{Advances in Neural Information Processing Systems}, 2024.

\bibitem{ke2024era3d}
Ke Q, et al. Era3D: High-resolution multiview diffusion using efficient row-wise attention. arXiv:2405.11616, 2024.


\bibitem{shao2025mvpainter}
Shao M, Xiong F, Sun Z, Xu M. MVPainter: Accurate and detailed 3D texture generation via multi-view diffusion with geometric control. arXiv:2505.12635, 2025.

\bibitem{huang2024mvadapter}
Huang Z, et al. MV-Adapter: Multi-view consistent image generation made easy. arXiv:2412.03632, 2024.

\bibitem{chen2023text2tex}
Chen D Z, et al. Text2Tex: Text-driven texture synthesis via diffusion models. arXiv preprint, 2023.

\bibitem{richardson2023texture}
Richardson E, et al. TEXTure: Semantic texture transfer using diffusion models. arXiv preprint, 2023.

\bibitem{poole2023dreamfusion}
Poole B, Jain A, Barron J T, Mildenhall B. DreamFusion: Text-to-3D using 2D diffusion. \textit{International Conference on Learning Representations}, 2023.

\bibitem{chen2023fantasia3d}
Chen R, Chen Y, Jiao N, Jia K. Fantasia3D: Disentangling geometry and appearance for high-quality text-to-3D content creation. \textit{Proceedings of the IEEE/CVF International Conference on Computer Vision}, 2023.

\bibitem{zhang2023controlnet}
Zhang L, Rao A, Agrawala M. Adding conditional control to text-to-image diffusion models. \textit{Proceedings of the IEEE/CVF International Conference on Computer Vision}, 2023.

\bibitem{ye2023ipadapter}
Ye H, Zhang J, Liu S, Han X, Yang W. IP-Adapter: Text compatible image prompt adapter for text-to-image diffusion models. arXiv:2308.06721, 2023.


\bibitem{hu2022lora}
Hu E J, Shen Y, Wallis P, et al. LoRA: Low-rank adaptation of large language models. \textit{International Conference on Learning Representations}, 2022.

\bibitem{liu2024dora}
Liu S Y, Wang C Y, Yin H, et al. DoRA: Weight-decomposed low-rank adaptation. arXiv:2402.09353, 2024.

\bibitem{zhang2023adalora}
Zhang Q, Zuo S, Liang C, et al. AdaLoRA: Adaptive budget allocation for parameter-efficient fine-tuning. \textit{International Conference on Learning Representations}, 2023.

\bibitem{ho2020ddpm}
Ho J, Jain A, Abbeel P. Denoising diffusion probabilistic models. \textit{Advances in Neural Information Processing Systems}, 2020.

\bibitem{song2021ddim}
Song J, Meng C, Ermon S. Denoising diffusion implicit models. \textit{International Conference on Learning Representations}, 2021.

\bibitem{ho2022cfg}
Ho J, Salimans T. Classifier-free diffusion guidance. \textit{NeurIPS Workshop}, 2022.

\bibitem{castillo2023adaptive}
Castillo A, Legenstein R, et al. Adaptive guidance: Training-free acceleration of conditional diffusion models. arXiv:2312.12487, 2023.

\bibitem{wang2004ssim}
Wang Z, Bovik A C, Sheikh H R, Simoncelli E P. Image quality assessment: From error visibility to structural similarity. \textit{IEEE Transactions on Image Processing}, 2004, 13(4): 600--612.


\bibitem{zhang2018lpips}
Zhang R, Isola P, Efros A A, Shechtman E, Wang O. The unreasonable effectiveness of deep features as a perceptual metric. \textit{Proceedings of the IEEE/CVF Conference on Computer Vision and Pattern Recognition}, 2018.

\bibitem{radford2021clip}
Radford A, Kim J W, Hallacy C, et al. Learning transferable visual models from natural language supervision. \textit{International Conference on Machine Learning}, 2021.

\bibitem{oquab2024dinov2}
Oquab M, Darcet T, Moutakanni T, et al. DINOv2: Learning robust visual features without supervision. \textit{Transactions on Machine Learning Research}, 2024.

\bibitem{rombach2022ldm}
Rombach R, Blattmann A, Lorenz D, Esser P, Ommer B. High-resolution image synthesis with latent diffusion models. \textit{Proceedings of the IEEE/CVF Conference on Computer Vision and Pattern Recognition}, 2022.

\bibitem{esser2024scaling}
Esser P, Kulal S, Blattmann A, et al. Scaling rectified flow transformers for high-resolution image synthesis. \textit{International Conference on Machine Learning}, 2024.

\bibitem{blackforestlabs2024flux}
Black Forest Labs. FLUX: Flow matching for image generation. 2024.

\bibitem{loshchilov2019decoupled}
Loshchilov I, Hutter F. Decoupled weight decay regularization. \textit{International Conference on Learning Representations}, 2019.

\bibitem{cherti2023reproducible}
Cherti M, Beaumont R, Wightman R, et al. Reproducible scaling laws for contrastive language-image learning. \textit{Proceedings of the IEEE/CVF Conference on Computer Vision and Pattern Recognition}, 2023.

\bibitem{nichol2022glide}
Nichol A, Dhariwal P, Ramesh A, et al. GLIDE: Towards photorealistic image generation and editing with text-guided diffusion models. \textit{International Conference on Machine Learning}, 2022.

\bibitem{saharia2022photorealistic}
Saharia C, Chan W, Saxena S, et al. Photorealistic text-to-image diffusion models with deep language understanding. \textit{Advances in Neural Information Processing Systems}, 2022.

\bibitem{peebles2023scalable}
Peebles W, Xie S. Scalable diffusion models with transformers. \textit{Proceedings of the IEEE/CVF International Conference on Computer Vision}, 2023.

\end{thebibliography}
\end{document}
